\documentclass[3p,times,authoryear]{elsarticle}

% ---------------------------------------------------------------------------
% Prepared for submission to Ecological Modelling (Elsevier format)
% ---------------------------------------------------------------------------

\usepackage[utf8]{inputenc}
\usepackage[T1]{fontenc}
\usepackage{lmodern}
\usepackage{microtype}
\usepackage{amsmath,amssymb}
\usepackage{booktabs}
\usepackage[hidelinks]{hyperref}

\newcommand{\gssk}{\textsc{gssk}}
\newcommand{\esl}{\textsc{esl}}
\newcommand{\idc}{\textsc{idc}}
\newcommand{\emsim}{\textsc{EmSim}}
\newcommand{\latexenergese}{\texttt{latex-energese}}
\newcommand{\code}[1]{\texttt{\small #1}}

\journal{Ecological Modelling}

\begin{document}

\begin{frontmatter}

\title{Architectural precedents in systems ecology modelling: A retrospective analysis of the Emergy Simulator (\emsim) and implications for the Energese project}

\author{Sholto Maud}
\affiliation{organization={Energese Project}, city={Sydney}, country={Australia}}

\begin{abstract}
The computational formalisation of Odum's Energy Systems Language (\esl) remains a persistent challenge in ecological modelling. In 2004, \citet{valyi2004} introduced \emsim, an open-source Java implementation designed to translate \esl{} diagrams into ordinary differential equations (\textsc{ode}s), integrate them, and propagate emergy across the resulting network. Two decades later, the Energese project (\gssk) attempts to solve the same fundamental parametric and topological problems. This paper presents a comparative architectural analysis of \emsim{} and \gssk. We evaluate the transition from embedded expression parsers to closed-vocabulary schemas, the algorithmic trade-offs between path enumeration and linear solvers for emergy propagation, and the shift in focus from file exchange to bit-identical execution reproducibility. The analysis demonstrates that while contemporary frameworks benefit from modern data schemas and deterministic runtimes, \emsim's 2004 implementation of Odum's co-product allocation rules via path enumeration remains a critical benchmark that current linear-solve architectures have yet to fully replicate. This retrospective aims to ground future \esl{} software development in established computational precedents.
\end{abstract}

\begin{keyword}
Energy Systems Language \sep Systems Ecology \sep Emergy \sep Ecological Modelling \sep Software Architecture
\end{keyword}

\end{frontmatter}

\section{Introduction}
\label{sec:intro}

The standardisation of systems ecology models requires robust computational frameworks capable of translating qualitative diagrams into quantitative, reproducible simulations. H.T. Odum's Energy Systems Language (\esl) provides a rigorous thermodynamic and parametric vocabulary for this purpose \citep{odum1994,odum1996}. However, the translation of \esl{} diagrams into executable ordinary differential equations (\textsc{ode}s) has historically been ad hoc, often relying on commercial software that obscures the underlying network algebra \citep{valyi2004}.

The Emergy Simulator (\emsim), developed by \citet{valyi2004}, represents one of the earliest comprehensive attempts to build a dedicated, open-source computational platform for \esl. \emsim{} was designed to parse diagrams, generate \textsc{ode}s, execute numerical integration, and calculate emergy transformities dynamically. Currently, the Energese project (comprising the \gssk{} C99 kernel, a LuaLaTeX drawing package, and a browser-based editor) is developing a contemporary infrastructure with identical core objectives. 

This paper analyses the \emsim{} specification as a design precedent. By comparing the algorithmic and architectural decisions documented by \citet{valyi2004} against the current \gssk{} implementation, we identify persistent challenges in ecological systems modelling, specifically concerning network algebra, numerical integration, and execution reproducibility.

\section{Core Motivations for Computational \esl}
\label{sec:motivations}

\citet{valyi2004} articulated six primary motivations for \emsim, which remain highly relevant to contemporary systems ecology:
\begin{enumerate}
    \item \textbf{Standardisation:} Models require open implementations and specific file formats to be shareable and subject to peer review.
    \item \textbf{Independence:} Modelling should not depend on proprietary commercial software.
    \item \textbf{Deterministic Translation:} The translation of diagrams to \textsc{ode}s must follow strict, automated rules rather than ad hoc transcription.
    \item \textbf{Dynamic Transformity:} Transformities must be computed within the specific system context rather than imported from static tables, which risks double-counting and invalidates emergy indicators.
    \item \textbf{Extensibility:} Software must be modular and open-source to support cooperative academic development.
    \item \textbf{Portability:} The platform must run across different operating systems.
\end{enumerate}

The Energese project shares these foundational requirements. However, the computational methods used to achieve them have evolved. The following sections detail the architectural divergences between the two approaches.

\section{Comparative Architectural Analysis}
\label{sec:architecture}

\subsection{Data Structures and Persistence}
\emsim{} utilised an XML formalism, serialising the Java object graph (JavaBeans) to store view attributes, graph topology, and mathematical models (including behavioural methods) within each node \citep{valyi2004}. While XML provided human readability, coupling the file format directly to the Java class graph created fragility; the semantics of the model were dependent on the presence of specific Java classes at runtime.

In contrast, \gssk{} separates the model definition from the runtime environment using a strict JSON Schema (\code{gssk.schema.json}). Behaviour is not stored as executable code within the file. Instead, the schema defines a closed vocabulary of declarations (e.g., nine primitive node types). This ensures that the model document acts as a pure parametric specification, independent of the C99 kernel that executes it.

\subsection{Mathematical Evaluation and Parsing}
To evaluate user-defined expressions, \emsim{} embedded the FrAid solver, allowing arbitrary mathematical formulae (e.g., polynomials, trigonometry) \citep{valyi2004}. While highly flexible, arbitrary string parsing introduces significant limitations: it inhibits symbolic analysis (such as automated Jacobian derivation or conservation checks) and introduces floating-point evaluation ambiguities across different hardware.

\gssk{} replaces the arbitrary parser with closed enumerations: eight edge logics (e.g., \code{constant}, \code{linear}, \code{interaction}, \code{reversible}) and eight forcing waveforms. While this restricts the modeller to a predefined mathematical vocabulary, it guarantees that every model in the corpus is analytically tractable. Furthermore, it allows the kernel to identify linear states and apply exact integration methods, a guarantee that cannot be made when parsing arbitrary strings.

\subsection{Emergy Propagation and Network Algebra}
The most significant algorithmic divergence between \emsim{} and \gssk{} lies in the propagation of emergy across complex networks, specifically regarding Odum's rules for co-production. 

Odum stipulates that when a process yields multiple by-products, each by-product carries the total emergy of the inputs; however, if these by-products later recombine, their emergy must not be double-counted \citep{odum1996, brown1996}. \citet{valyi2004} implemented this by enumerating all pathways from a source to a target. The algorithm sorted pathways into co-produced sets and retained only the maximum emergy contribution from each set, effectively preventing double-counting. 

While mathematically faithful to Odum's rules, path enumeration scales exponentially ($O(2^n)$) in dense networks, leading to severe performance bottlenecks. Consequently, the field largely shifted toward matrix and eigenvector approaches \citep{collins2000}. 

\gssk{} currently utilises a linear solve for quality propagation. While computationally efficient and exact for flow splits, this approach currently fails to implement the bounding rule for co-products. If $n$ co-product pathways recombine at a node in \gssk, the node incorrectly reports the sum of the copies rather than the maximum. Therefore, \emsim's 2004 implementation remains algorithmically superior in its adherence to \esl{} rules, highlighting a critical area for future \gssk{} development. A proposed solution involves a longest-path computation in the $(\max,+)$ semiring per source, which would preserve provenance without the exponential cost of full path enumeration.

\subsection{Discontinuities and Numerical Integration}
\emsim{} handled discontinuous variations using a switch component that generated an ``almost infinite'' slope, integrated via a fixed-step 4th-order Runge-Kutta (RK4) method \citep{valyi2004}. This approach introduces stiffness; fixed-step explicit integrators struggle with near-infinite slopes, resulting in trajectory errors depending on where the step boundaries fall.

\gssk{} implements strict \code{switch} nodes and \code{threshold} logics. If a condition flips within a time step (\code{dt}), the integration proceeds as though the state remained on its initial side until the next boundary. While this bounds the residual error and makes it predictable, \gssk{} currently lacks event detection or step-boundary refinement. Both systems highlight the inherent difficulty of modelling discrete events in continuous ecological simulations.

\section{Reproducibility and Software Longevity}
\label{sec:reproducibility}

A critical omission in the 2004 \emsim{} specification is the concept of execution reproducibility. \citet{valyi2004} focused extensively on sharing model files, but did not address determinism, random seed capture, or floating-point consistency. In computational modelling, a shared file is insufficient if the execution environment cannot be exactly replicated to audit the resulting emergy indicators.

\gssk{} treats bit-identical reproduction as a baseline requirement. Nondeterminism is strictly avoided, structural mutations are logged, and stochastic runs draw from an instance-owned SplitMix64 generator with captured seeds and stream positions. 

Furthermore, the decay of \emsim's software dependencies (Java applets, specific JGraphpad versions) demonstrates that software applications are ephemeral. \gssk's reliance on C99 and WebAssembly is an attempt to target stable, low-level standards rather than vendor-specific runtimes. The archival unit of ecological modelling must be the model document and its schema specification, not the executable program.

\section{Discussion and Future Directions}
\label{sec:discussion}

\citet{valyi2004} noted suggestions by \citet{giannantoni2002} regarding the necessity of accumulation terms and fractional calculus for dynamic empower studies. \emsim{} lacked the solvers to implement this. \gssk{} has since implemented the Incipient Differential Calculus (\idc) \citep{giannantoni2006}, providing exact closed-form integration for isolated interactions and Padé-(3,3) matrix exponentials for cyclic networks. However, the theoretical debate regarding accumulation terms remains open, and \gssk{} records discrepancies between classical RK4 and \idc{} as observable confidence flags rather than resolved scientific consensus.

Moving forward, the Energese project must address the co-product double-counting defect in its linear solver. Additionally, while \emsim{} proposed embedding statistical distributions and confidence intervals directly into the model data structure for stochastic simulation, \gssk{} currently handles stochasticity externally via ensemble forecasting. Integrating parameter uncertainty directly into the \gssk{} JSON schema represents a logical next step.

\section{Conclusion}
\label{sec:conclusion}

The 2004 \emsim{} project accurately identified the primary computational bottlenecks in systems ecology: the need for deterministic \textsc{ode} translation, dynamic transformity calculation, and the prevention of double-counting in co-production networks. While modern frameworks like \gssk{} offer superior data schemas, deterministic execution, and advanced calculus implementations, \emsim's path-enumeration approach to Odum's co-product rules remains a vital benchmark. By treating historical software implementations as rigorous design precedents, contemporary ecological modelling can avoid repeating past architectural compromises and build more robust, auditable tools for sustainable development analysis.

\begin{thebibliography}{7}

\bibitem[{Brown and Herendeen(1996)}]{brown1996}
Brown, M.\,T., Herendeen, R.\,A., 1996. Embodied energy analysis and emergy analysis: a comparative view. Ecological Economics 19(3), 219--235.

\bibitem[{Collins and Odum(2000)}]{collins2000}
Collins, D., Odum, H.\,T., 2000. Calculating Transformities with an Eigenvector Method. Center for Environmental Policy, University of Florida, Gainesville, FL.

\bibitem[{Giannantoni(2002)}]{giannantoni2002}
Giannantoni, C., 2002. The Maximum Em-Power Principle as the Basis for Thermodynamics of Quality. Servizi Grafici Editoriali, Padua.

\bibitem[{Giannantoni(2006)}]{giannantoni2006}
Giannantoni, C., 2006. Mathematics for generative processes: living and non-living systems. Journal of Computational and Applied Mathematics 189(1--2), 324--340.

\bibitem[{Odum(1994)}]{odum1994}
Odum, H.\,T., 1994. Ecological and General Systems: An Introduction to Systems Ecology. University Press of Colorado, Niwot.

\bibitem[{Odum(1996)}]{odum1996}
Odum, H.\,T., 1996. Environmental Accounting: Emergy and Environmental Decision Making. John Wiley \& Sons, New York.

\bibitem[{Valyi and Ortega(2004)}]{valyi2004}
Valyi, R., Ortega, E., 2004. Emergy simulator, an open source simulation platform dedicated to systems ecology and emergy studies. In: Ortega, E., Ulgiati, S. (Eds.), Proceedings of the IV Biennial International Workshop ``Advances in Energy Studies''. Unicamp, Campinas, SP, Brazil, pp. 349--360.

\end{thebibliography}

\end{document}
